\documentclass[10pt,a4paper]{amsart}
\usepackage[margin=19mm]{geometry}
\usepackage{amsmath,amssymb,mathtools}
\usepackage[scr=rsfs,cal=euler]{mathalfa}
\usepackage{xfrac}
\usepackage[unicode,hidelinks,bookmarks=false]{hyperref}
\newcommand{\ie}{i.e.\ }
\newcommand{\cf}{cf.\ }
\newcommand{\resp}{respectively\ }
\newcommand{\Subsection}{\subsection}
\providecommand{\nobreakdash}{\nobreak\hbox{-}}
\setlength{\emergencystretch}{2em}
\multlinegap0pt
\usepackage{amssymb}
\usepackage{bm}
\usepackage[shortlabels]{enumitem}
\newtheorem{thm}{Theorem}
\newtheorem{lem}{Lemma}
\newtheorem{proposition}{Proposition}
\renewcommand\ge{\geqslant}
\renewcommand\le{\leqslant}
\newcommand{\remove}[1]{}
\title{Bounds for Characters of the Symmetric Group: A Hypercontractive Approach}
\author{Noam Lifshitz, Avichai Marmor}
\date{Source: arXiv:2308.08694v4 (CC BY 4.0)}
\begin{document}
\maketitle

\noindent\emph{Source and licence.} Bounds for Characters of the Symmetric Group: A Hypercontractive Approach. Version arXiv:2308.08694v4 (CC BY 4.0), distributed by arXiv under the Creative Commons Attribution 4.0 International licence, http://creativecommons.org/licenses/by/4.0/, as recorded in the arXiv licence metadata for this exact version and shown on the abstract page https://arxiv.org/abs/2308.08694v4. Full licence notice: \texttt{licenses/arxiv-2308.08694-CC-BY-4.0.txt}.\par\smallskip

\noindent\textit{Editorial scope.} Counted unit: the article's thm main (source0.tex lines 378--383). The article's complete original proof of it is delivered with it (source0.tex lines 835--878, one \texttt{proof} environment of 44 lines, which the article places away from the statement). No mathematics is written, repaired or re-derived: the statement and the proof are the article's own lines, in the article's own notation, and no retained line is substituted or re-flowed.  Retained uncounted as the source-local context the proof needs: lines 584--587; lines 593--597; lines 752--754; lines 822--824.  Nothing was omitted from any retained span, so the package carries no omission record.  No retained line contains a citation, so the wrapper carries no bibliography and no bibliography entry.  No retained line contains a figure environment, an image-inclusion command or drawing code, so no image is delivered and the package declares no asset dependency.  Editorial formatting only, in addition: an AMS \texttt{amsart} preamble that restates the article's own declarations and macros used by the retained lines, namely the environments \texttt{thm}, \texttt{lem}, \texttt{proposition} on separate counters, together with the standard packages those lines need.  The retained environments print in this document as Theorem 1, Lemma 1, Proposition 1 in order of appearance, whereas the article prints the same items with the numbering of its own full document.  The complete native source of the article as distributed in the arXiv e-print for this exact version is bundled unchanged as \texttt{sources/145742/source0.tex}.
\begin{lem}\label{lem: dimension of characters}
Let $d$ be a positive integer and let $\lambda\vdash n$ of level $d$. Denote $\chi := \chi_\lambda$ and $\tilde{\chi} := \chi_{\tilde{\lambda}}$.
Then $\dim(\chi) \ge \binom{n-d}{d}\dim(\tilde{\chi})$.
\end{lem}

\begin{lem}\label{lem: dimension of characters converse}
Let $d$ be an integer and let $\lambda\vdash n$ of level $d$. Denote $\chi := \chi_\lambda$ and $\tilde{\chi} := \chi_{\tilde{\lambda}}$.
Then 
\[\dim(\chi) \le \binom{n}{d}\dim(\tilde{\chi}).\]
\end{lem}

\begin{lem} \label{lem:murnaghan nakayama remove long cycles}
    Let $\lambda \vdash n$ be a partition with $\lambda_1 = n-d$, and let $\ell \ge d + \lambda_2$. Let $\mu = (\lambda_1 - n + \ell, \lambda_2, \ldots, \lambda_r)$ be the partition obtained from $\lambda$ by removing $n-\ell$ boxes from its first row. Let $\sigma \in S_n$ be a permutation with a cycle structure $1^\ell \alpha_1\cdots \alpha_k$ with $\alpha_i > d$. That is, $\sigma$ has $\ell$ fixed points, and the rest of the cycles of $\sigma$ have length at least $d+1$. Then we have $\chi_{\lambda}(\sigma) = \chi_{\mu}(1)$.
\end{lem}

\begin{proposition} \label{prop:probability no short cycles}
    Let $d$ and let $r>d$. Let $F_{r,d}$ be the event that a random permutation $\sigma\sim S_r$ has all its cycles of length $>d$. Then $\Pr[F_{r,d}]\ge \frac{1}{10d}$.
\end{proposition}

\begin{thm}\label{thm:matching lower bound for main}
There exists an absolute constant $c>0$, such that the following holds. Let $q\ge 2$ and let $d < \min(e^{cq},\frac{n}{q+1})$ be an integer. Let $\chi$ be a character of level $d$. Then  
\[
\| \chi\|_q \ge \left( \frac{cq}{\log (qd)} \right)^d\left(\frac{d^d \chi(1)}{n^d}\right)^{1-2/q}. 
\]
\end{thm}

\begin{proof}[Proof of Theorem \ref{thm:matching lower bound for main}]
Let $d < \min(e^{cq},\frac{n}{q+1})$ and let $\chi_{\lambda}$ be a character of level $d$. As multiplication by sign does not affect the $q$-norm we may assume without loss of generality that $\lambda_1 = n-d$. Pick some positive integer $\ell \ge 2d$ that will be determined later. Let $\sigma$ be a permutation with $\ell$ fixed points such that rest of the cycles of $\sigma$ have length at least $d+1$. Let $\mu = (\lambda_1 - n + \ell, \lambda_2, \ldots, \lambda_r)$ be obtained from $\lambda$ by removing $n-\ell$ boxes from its first row. As usual, let $\tilde{\lambda}$ be obtained from $\lambda$ by deleting its first row. Since $\lambda_2 \le d$, we may apply Lemma \ref{lem:murnaghan nakayama remove long cycles} and obtain $\chi_{\lambda}(\sigma) = \chi_{\mu}(1)$.
By Lemmas \ref{lem: dimension of characters} and \ref{lem: dimension of characters converse} we have
\[\chi_{\mu}(1)\ge \binom{\ell-d}{d} \chi_{\tilde{\lambda}}(1)\ge \frac{\binom{\ell-d}{d}}{\binom{n}{d}} \chi_{\lambda}(1)\ge \left(\frac{c\ell}{n}\right)^d \chi_{\lambda}(1)\]
for an absolute constant $c$. Let $E$ be the event that a random $\sigma\sim S_n$ has $\ell$ fixed points and all other cycles of length $>d$. Then we may %apply Lemma \ref{lem: dimension of characters converse} to
obtain that
\[
\| \chi_{\lambda} \|_q^q \ge \Pr[E] \chi^q_\mu(1) \ge \Pr[E] \left(\frac{c\ell}{n}\right)^{dq}\chi^q_{\lambda}(1). 
\]

Let $F_{r,d}$ be the event that a random permutation $\sigma\sim S_r$ has all its cycles of length $>d$. Proposition~\ref{prop:probability no short cycles} implies that whenever $n-\ell > d$, we have $\Pr[F_{n-\ell,d}]\ge \frac{1}{10d}$. Therefore,
% \[ \Pr[E] = \binom{n}{\ell}\frac{1}{n(n-1)\cdots 
% (n-\ell +1)}\Pr[F_{n-\ell,d}] = \frac{1}{\ell!}\Pr[F_{n-\ell,d}] \ge \frac{1}{10d \ell!} \ge \frac{1}{ (c\ell)^\ell d}\]
\begin{eqnarray*}
\Pr[E] &=& \binom{n}{\ell}\frac{1}{n(n-1)\cdots (n-\ell +1)}\Pr[F_{n-\ell,d}] \\
       &=& \frac{1}{\ell!}\Pr[F_{n-\ell,d}] \ge \frac{1}{10d \ell!} \ge \frac{1}{ (c\ell)^\ell d}
\end{eqnarray*}
for an absolute constant $c$.

\remove{It therefore remains to lower bound $\Pr[F_{n-l}]$. %Now the cycle decomposition of a random permutation in $F_r$ has the distribution obtained by first choosing a random $i_1\sim r$ and then repeating with $r-i_1$ in place of $r$. This follows easily from the fact that $r$ participates in a cycle of random length and that when removing its cycle we get a uniformly random permutation. 
The distribution of the cycle decomposition of a random permutation in $S_r$ is identical to the distribution of a random sequence $(i_1,i_2,\dots)$ obtained by the following process: First choose a random $i_1\sim [r]$, then repeat the process recursively with $r-i_1$ in place of $r$ to obtain $(i_2,\dots)$. This follows easily from the fact that the length of the cycle of the element $r$ is uniform, and that when removing its cycle we get a uniformly random permutation on the remaining elements. 

This yields the recursive formula \[ \Pr[F_r] = \frac{1}{r} + \sum_{i=d+1}^{r-d-1}\frac{1}{r}\Pr[F_{r-i}].\]
We now show that $\Pr[F_r]\ge \frac{1}{10d}$ by induction on $d$ for all $r> d$. For $r \in (d+1,10 d)$ the claim is obvious, since $\Pr[F_r] \ge \frac{1}{r}$ by the recursive formula. We now apply the induction hypothesis to obtain that %\[\Pr[F_r]\ge \frac{1}{n}+\frac{n-11 d}{10 d n} + \sum_{i=d+1}^{10d} \frac{1}{ni} > \frac{1}{10 d }-\frac{11}{10n} + \frac{\log(10 d) - \log d}{n}>\frac{1}{10 d}.\] % a little inaccurate and complicated
\[
    \Pr[F_r]\ge \frac{1}{r} + \sum_{i=d+1}^{r-d-1} \frac{1}{10dr} = \frac{1}{r}+\frac{r-2d-1}{10dr}>\frac{1}{10 d}.
\]

This shows that whenever $n-\ell >d$, we have 
\[\Pr[E]\ge \frac{1}{10d \ell!} \ge \frac{1}{ (c\ell)^\ell d}\]
for an absolute constant $c$.}
Substituting $\ell=\frac{qd}{\log (qd)},$ while noting that $\ell\ge 2d$ and $n - \ell > d$ provided that $d < \min(e^{cq},\frac{n}{q+1})$, we obtain  
% \[
% \|\chi_{\lambda}\|_q \ge \left(\frac{c^2qd}{\log (qd)}\right)^d\frac{\chi_\lambda(1)}{n^{d}} =  \left(\frac{c^2q}{\log(qd)}\right)^d\frac{d^d\chi_{\lambda}(1)}{n^d}\ge \left(\frac{c'q}{\log(qd)}\right)^d \left(\frac{d^d\chi_{\lambda}(1)}{n^d}\right)^{1-2/q},
% \]
\begin{eqnarray*}
\|\chi_{\lambda}\|_q 
&\ge& \left(\frac{c^2qd}{\log (qd)}\right)^d\frac{\chi_\lambda(1)}{n^{d}} \\
&=& \left(\frac{c^2q}{\log(qd)}\right)^d\frac{d^d\chi_{\lambda}(1)}{n^d} \\
&\ge& \left(\frac{c'q}{\log(qd)}\right)^d \left(\frac{d^d\chi_{\lambda}(1)}{n^d}\right)^{1-2/q}
\end{eqnarray*}
for an absolute constant $c'$. The last inequality follows from the fact that $\chi_\lambda(1) \ge \binom{n-d}{d},$ which yields
$\left(\left(\frac{d}{n}\right)^d\chi_{\lambda}(1)\right)^{2/q} \ge \left(\frac{1}{3}\right)^{d}$.
\end{proof}

\end{document}
